\documentclass[11pt]{article}

\usepackage[margin=1in]{geometry}
\usepackage{amsmath,amssymb}
\usepackage{mathtools}
\usepackage{booktabs}
\usepackage{xcolor}
\usepackage{enumitem}
\usepackage[hidelinks,colorlinks=true,linkcolor=blue!50!black,urlcolor=blue!60!black]{hyperref}
\usepackage{fancyhdr}
\usepackage{parskip}
\usepackage{titlesec}
\usepackage{array}

\titleformat{\section}{\large\bfseries}{\thesection.}{0.5em}{}
\titleformat{\subsection}{\normalsize\bfseries}{\thesubsection}{0.5em}{}

\pagestyle{fancy}
\fancyhf{}
\fancyhead[L]{COSC 81/281, Spring 2026}
\fancyhead[R]{Final Project Proposal (FP2)}
\fancyfoot[C]{\thepage}

\title{\textbf{Learned Map-Completion Priors for Mobile-Robot Frontier Exploration}\\
\large Final Project Proposal (FP2)}
\author{Moin Mattar \\ COSC 81/281, Spring 2026}
\date{2026-05-20}

\begin{document}
\maketitle
\thispagestyle{fancy}

\noindent\textbf{Format:} Solo $\cdot$ ROS 2 Humble $\cdot$ Stage simulator $\cdot$ builds directly on PA3 + PA4.\\
\textbf{Track:} Frontier Explorer (per the spec's example tracks).

\textit{Integrity \& AI use. This work is my own. AI helped me with LaTeX formatting and literature search. Collaborators: none. Signature: Moin Mattar.}

\section*{Abstract}
Frontier-based exploration scores each candidate frontier with a hand-crafted heuristic that is blind to the structure of the unseen environment. We propose to learn a generative prior over completed occupancy grids using a small \textbf{conditional denoising diffusion model}, and use samples from that prior to score frontiers by \emph{expected} information gain with an \emph{uncertainty bonus}. The method is trained on synthetic (partial-map, full-map) pairs from the \textbf{HouseExpo} dataset and evaluated in Stage on the maze world from PA3 plus additional procedurally generated worlds. We hypothesize that sampling-based scoring reduces time-to-coverage relative to standard frontier exploration, with the largest gains in environments with structural ambiguity (corridor vs. dead-end). We discuss extensions to \textbf{Diffusion Forcing} for variable-horizon rollouts as future work.

\section{Introduction}

Standard frontier-based exploration (Yamauchi 1997; Quin et al. 2014, the algorithm covered in Lec 19) picks the next frontier with a hand-crafted score, typically:
\[
\text{score}(f) = \alpha \cdot N_\text{unknown}(f) - \beta \cdot d(r, f),
\]
where $N_\text{unknown}(f)$ is the number of unknown cells visible from frontier $f$ and $d(r,f)$ is the robot's distance to it. This heuristic ignores \textbf{what lies beyond the frontier}: a frontier opening into a wide unmapped room and a frontier opening into a dead-end closet can score identically; only after the robot drives there does the difference manifest.

Learned map-completion priors address this. Shrestha et al. (ICRA 2019) trained a conditional GAN to predict the full map from a partial observation and used the predicted free-space to inform frontier scoring. More recent work uses diffusion models (Lin et al. 2024, arXiv:2409.10681) for 3D occupancy prediction at the frontier. Both these methods produce a \textbf{single deterministic prediction} and use it as a point estimate of ``what's behind the wall.''

Our contribution is to treat the unseen environment as \textbf{inherently uncertain} and score frontiers by the \emph{distribution} of plausible completions. For each candidate frontier $f$, we sample $K$ completions from the diffusion model, then score $f$ by $\mathbb{E}[\text{info gain}] + \lambda \cdot \text{Var}[\text{info gain}]$. The variance bonus encourages the robot toward frontiers where the model is \emph{uncertain}, which are precisely the frontiers where new observations are most informative.

\section{Method}

\subsection{Map Completion Model}

\textbf{Architecture.} A small conditional U-Net ($\sim$5M parameters) operating on $256\times256$ occupancy grid patches. Inputs: a partial grid where each cell is one of \{free, occupied, unknown\} (encoded as a 3-channel one-hot image). Output: a predicted full grid (3 channels, softmaxed per cell).

\textbf{Training objective.} Standard conditional DDPM (Ho et al. 2020):
\[
\mathcal{L} = \mathbb{E}_{t, x_0, \epsilon}\left[\| \epsilon - \epsilon_\theta(x_t, t, c) \|^2\right],
\]
where $c$ is the partial map (conditioning) and $x_0$ is the ground-truth full map.

\textbf{Training data.} \textbf{HouseExpo} dataset (Li et al. IROS 2020): 35,000+ indoor floor plans. For each map we generate (partial, full) pairs by simulating a partial-observation mask: place a virtual robot at a random cell, cast a 360$^\circ$ lidar with realistic range and angular resolution (matching the Stage rosbot's lidar), mark observed cells as known, the rest as unknown. The resulting partial map is the conditioning $c$, the original full map is $x_0$. This yields $\sim$500k training samples cheaply.

\subsection{Sampling-Based Frontier Scoring}

For each candidate frontier $f$ at inference time:
\begin{enumerate}[itemsep=0pt,topsep=0pt]
  \item Crop the current partial map to a $256\times256$ window centered on the robot.
  \item Sample $K = 8$ completions from the diffusion model conditioned on the partial map.
  \item For each completion, compute the realized info gain: the number of currently-unknown cells that the completion predicts as \textbf{free and reachable from $f$} (within a configurable radius).
  \item Score $f$ by
  \[
  \text{score}(f) = \mathbb{E}_k[g_k(f)] + \lambda \cdot \text{Std}_k[g_k(f)] - \beta \cdot d(r, f),
  \]
  where $g_k(f)$ is the info gain under completion $k$, $\lambda$ controls the uncertainty bonus, and the last term is the standard distance penalty.
  \item Pick $f^* = \arg\max_f \text{score}(f)$, drive to it, scan, update the partial map, repeat.
\end{enumerate}

\subsection{ROS 2 Integration}

The system is one ROS 2 node that subscribes to \texttt{/map} (from a slightly modified PA4 mapper that also tracks frontier cells) and publishes both \texttt{/predicted\_map} (the mean completion, for visualization) and \texttt{/scored\_frontiers} (a \texttt{geometry\_msgs/PoseArray}, each pose annotated with its score). A simple pure-pursuit controller drives the robot to the highest-scoring frontier; the loop continues.

\section{Connection to Course Material}

\begin{tabular}{p{0.32\textwidth} p{0.6\textwidth}}
\toprule
\textbf{Class component} & \textbf{Where in this project} \\
\midrule
Lec 10, occupancy grids & The map representation; \texttt{pa4\_grid\_mapper.py} is the substrate. \\
Lec 12, search-based planning & A* used to compute $d(r, f)$ and reachability inside each completion. \\
Lec 14--16, Bayes filtering & The PA4 log-odds update; the diffusion model is a learned generative version of the same posterior. \\
Lec 19, exploration \& coverage & The frontier-based exploration loop (Quin et al. 2014). \\
Lec 20, ML for robotics & The conditional diffusion model is the ML component. \\
\bottomrule
\end{tabular}

\medskip
\noindent The project synthesizes \textbf{five distinct lecture units} into one autonomous system.

\section{Experimental Plan}

\subsection{Worlds}
PA3 maze (small, structured), two procedurally generated mazes (one with dead-ends, one with one large central room), and two HouseExpo floorplans loaded into Stage (out-of-training-distribution). \textbf{Five worlds $\times$ ten trials each = 50 trials per condition.}

\subsection{Baselines}
\begin{itemize}[itemsep=0pt,topsep=0pt]
  \item \textbf{Heuristic frontier.} Standard $\text{score}(f) = N_\text{unknown}(f) - \beta d(r, f)$.
  \item \textbf{Shrestha-style.} A deterministic conditional U-Net (no diffusion) trained on the same data, single prediction used for info gain. Isolates ``stochastic sampling'' from ``learned prior.''
  \item \textbf{Ours (no variance bonus).} $\lambda = 0$. Isolates ``expected info gain'' from ``uncertainty-seeking.''
  \item \textbf{Ours (full).} $\lambda > 0$.
\end{itemize}

\subsection{Metrics}
\begin{itemize}[itemsep=0pt,topsep=0pt]
  \item \textbf{Time-to-90\%-coverage} (primary).
  \item \textbf{Path length} to 90\% coverage.
  \item \textbf{Info-gain calibration}: reliability diagram of predicted vs. realized info gain.
  \item \textbf{Qualitative}: side-by-side visualization of $K=8$ completions per frontier (the headline demo figure).
\end{itemize}

\subsection{Ablations}
$K \in \{1, 4, 8, 16\}$ (sample-count sweep); $\lambda \in \{0, 0.5, 1.0, 2.0\}$ (variance-bonus sweep); training-data scale 10\%/50\%/100\%.

\section{Risks and Fallbacks}

\begin{tabular}{p{0.4\textwidth} p{0.55\textwidth}}
\toprule
\textbf{Risk} & \textbf{Mitigation} \\
\midrule
Diffusion training won't converge in time. & Diffusion U-Nets are well-trodden territory. Many working open-source implementations exist (\texttt{huggingface/diffusers}, \texttt{lucidrains/denoising-diffusion-pytorch}). If training fails, swap for a deterministic U-Net regressor: degrades to the Shrestha-style baseline but still novel (sampling becomes dropout-based). \\
\addlinespace
Sim-to-sim domain gap (HouseExpo $\to$ Stage). & Include 2 Stage worlds in training. Domain-randomize lidar noise. \\
\addlinespace
ROS 2 integration eats time. & The mapper, frontier extractor, and pure-pursuit controller are already implemented in PA3/PA4. Only the diffusion-inference module is new. \\
\addlinespace
Inference too slow for real-time. & At $256\times256$ and $K=8$, denoising takes $\sim$1--2s on a laptop GPU. Acceptable because frontier selection happens once every $\sim$5 seconds during driving. \\
\bottomrule
\end{tabular}

\section{Future Work, Diffusion Forcing for Variable-Horizon Rollouts}

The current method predicts a \emph{single-shot} completion. A natural extension is to predict the \textbf{evolution of the map over time} as the robot moves toward a candidate frontier, a video-generation problem rather than image-completion. \textbf{Diffusion Forcing} (Chen et al., NeurIPS 2024) is purpose-built for this: it trains a sequence model with independent per-token noise levels, enabling variable-horizon rollouts with calibrated per-step uncertainty. If time permits past FP3 (June 1), I will adapt the open-source \texttt{buoyancy99/diffusion-forcing} repo to roll out (map\_patch, action) sequences and score frontiers by expected info gain over a 20-step horizon. \textbf{This is flagged as stretch work, not a deliverable.}

\section{Timeline}

\begin{tabular}{l p{0.7\textwidth}}
\toprule
\textbf{Date} & \textbf{Milestone} \\
\midrule
May 20--22 & Repo setup. Reproduce a HuggingFace diffusion tutorial on MNIST to confirm pipeline. \\
May 22--26 & Data generation: HouseExpo loader, partial-map simulator, $\sim$500k training samples cached. \\
May 26--27 & First training run begins. \textbf{FP2.5 in-class progress update (May 27)}: show loss curves + sample completions. \\
May 27--30 & Training converges. Implement frontier scorer. Integrate into PA4 ROS node. \\
May 30--Jun 1 & Run baselines + ablations in Stage. Generate plots. \textbf{FP3 (June 1)}: 6-min presentation + demo video. \\
Jun 1--Jun 6 & Write 4-page IEEE paper. Polish demo. Polish github repo. \textbf{FP4 (June 6)}: final report + code + 3-min teaser. \\
\bottomrule
\end{tabular}

\section{Deliverables Summary}

4-page IEEE-format final report; open-source github repo with training + inference + ROS 2 integration; demo video showing the system in three Stage worlds; quantitative comparison against 3 baselines on 5 worlds $\times$ 10 trials.

\section{References}

\begin{itemize}[itemsep=0pt,topsep=0pt,leftmargin=*]
\item Yamauchi, B. \textit{A frontier-based approach for autonomous exploration.} CIRA 1997.
\item Quin, P. et al. \textit{Frontier-based multi-robot exploration with rendezvous.} 2014 (Lec 19).
\item Shrestha, R. et al. \textit{Learned map prediction for enhanced mobile robot exploration.} ICRA 2019.
\item Ho, J., Jain, A., Abbeel, P. \textit{Denoising diffusion probabilistic models.} NeurIPS 2020.
\item Li, T. et al. \textit{HouseExpo: a large-scale 2D indoor layout dataset for learning-based algorithms on mobile robots.} IROS 2020.
\item Lin et al. \textit{Online diffusion-based 3D occupancy prediction at the frontier with probabilistic map reconciliation.} 2024. arXiv:2409.10681.
\item Chen, B. et al. \textit{Diffusion forcing: next-token prediction meets full-sequence diffusion.} NeurIPS 2024. arXiv:2407.01392.
\end{itemize}

\vfill
\hrule
\medskip
\noindent\textbf{Talking points for the meeting (last page of the doc):}
\begin{itemize}[itemsep=0pt,topsep=0pt,leftmargin=*]
\item \textbf{One sentence:} learn a generative prior over completed occupancy grids; sample $K$ completions; score frontiers by expected info gain plus an uncertainty bonus.
\item \textbf{Why not Diffusion Forcing as the main deliverable:} adds complexity I can't safely deliver in three weeks solo. Single-shot version is cleaner to scope and evaluate; DF flagged as future work.
\item \textbf{Novelty over Shrestha 2019 / Lin 2024:} they predict one deterministic map and treat it as truth. We sample, average, and add a variance bonus, so the robot is aware of its own uncertainty.
\item \textbf{Ask the prof:} (1) does the novelty claim land; (2) Stage vs.\ Gazebo; (3) real robot at the end or sim-only; (4) is ``DF as future work'' the right framing.
\item \textbf{Insurance:} if diffusion training fails, the model degrades cleanly to a deterministic regressor. Sampling-based scoring still works (via dropout). Guaranteed minimum deliverable.
\end{itemize}

\end{document}
